% Corridor Attention --- a complete draft. Hard-wrapped at 80 columns.
\documentclass[11pt]{article}

\usepackage[utf8]{inputenc}
\usepackage{amsmath}
\usepackage{amssymb}
\usepackage{booktabs}
\usepackage{graphicx}
\usepackage{natbib}
\usepackage{amsthm}
\newtheorem{lemma}{Lemma}
\newtheorem{proposition}{Proposition}
\usepackage{hyperref}

\title{Corridor Attention: Bounded-Bandwidth Routing for Long-Context
Sequence Models}
\author{A.~Fixture \and B.~Fixture}
\date{}

\begin{document}
\maketitle

\begin{abstract}
Attention consumes a quadratic share of the compute in a transformer language
model, and that share sets a firm practical limit on how long a context one
can afford to train on. This paper presents \emph{Corridor Attention}, a
routing scheme in which each query token receives a fixed bandwidth of $k$ key
positions, taken from a learned, block-level corridor instead of from the
entire prefix. The layer retains a dense local window to handle short-range
syntax, adds a handful of summary tokens visible to every position to capture
discourse-level dependencies, and routes whatever bandwidth remains through a
differentiable block scorer trained jointly with the rest of the network.
What results is an attention layer whose cost grows linearly with sequence
length, with a constant the practitioner chooses directly. We present the
method, analyse its cost, and evaluate it on long-context language modelling
as well as on probes styled after retrieval tasks.
\end{abstract}

\section{Introduction}\label{sec:intro}

Self-attention \citep{vaswani2017attention} evaluates an interaction for each
pair of positions in a sequence. This is exactly why transformers capture
long-range structure so well, and equally why their cost scales as
$\Theta(n^2)$ with the sequence length $n$. At the context lengths that
production models are now routinely asked to handle --- tens of thousands of
tokens of source code, an entire scientific article together with its
appendices, a conversation spanning several days --- the attention term is no
longer a constant factor; it becomes the budget.

Previous work has responded along two broadly distinct lines. The first
family \emph{restricts} which positions may interact, using sliding windows,
dilated windows, strided or block-sparse masks, or hand-built combinations of
these three \citep{child2019sparse,beltagy2020longformer,zaheer2020bigbird}.
Such methods are inexpensive and predictable, yet their pattern is fixed
before the model has seen any data, so a dependency the designer failed to
foresee has simply no way to be expressed. The second family
\emph{approximates} the softmax kernel, putting a low-rank or feature-map
surrogate in place of the exact attention matrix
\citep{wang2020linformer,katharopoulos2020linear,choromanski2021performer}. They
retain a receptive field that is global in form, but the approximation error
lies outside the practitioner's control, and it interacts poorly with the
sharply peaked attention distributions that appear late in training.

Our own position is a third one, and we take it to be the pragmatic choice:
attention itself stays \emph{exact}, while \emph{the decision of what to
attend to} is learned. Every query is given a fixed bandwidth of $k$ key
positions, which we call a corridor. The blocks that the corridor passes
through are picked by a lightweight scorer that works on block summaries
instead of on individual positions. Since the bandwidth is set as a
hyperparameter and does not emerge from training, the layer's cost is known
before training begins; since the routing is learned, the pattern is free to
differ per head, per layer and per input.

This position commits us to three design choices, and it is these that set
Corridor Attention apart from the most closely related prior work.

\paragraph{Bandwidth is declared, not discovered.} It is the practitioner who
sets $k$, and the model can never go beyond it. Nothing in the training
objective trades accuracy for cost --- with the cost fixed, nothing is left to
trade. Capacity planning thereby becomes simple arithmetic instead of an
exercise in fitting empirical scaling curves.

\paragraph{Routing is scored on blocks, decided per query.} Were every
candidate key position to be scored, the quadratic term we set out to remove
would come straight back. Corridor Attention therefore scores
$\lceil n / b \rceil$ block summaries, so that the scorer costs
$\Theta(n^2 / b^2)$, a small fraction of the total at the block sizes we use.
The selection, however, is made per query, and two queries in one block may
be routed differently.

\paragraph{The local window is not negotiable.} Removing the dense local
window hurts perplexity far more than removing any equal amount of routed
bandwidth. For that reason the window is part of the architecture, not one
option among several, and the learned bandwidth goes only to what lies
beyond the window's reach.

We make three contributions. (i) We formalise routed attention under a
bounded bandwidth and give an instantiation of it, block-scored and
per-query, that trains end to end with no auxiliary load-balancing loss.
(ii) We derive a cost model for the layer and check it against measured
wall-clock time, finding close agreement between the predicted and the
observed crossover points against dense attention. (iii) We evaluate the
layer on long-context language modelling and on synthetic retrieval probes,
and we report the ablations that trace each benefit to one of the three
components.

\section{Background and Notation}\label{sec:background}

Consider an attention layer whose input is $X \in \mathbb{R}^{n \times d}$,
for $n$ positions and a model width of $d$. One head maps $X$ to queries,
keys and values through $Q = XW_Q$, $K = XW_K$ and $V = XW_V$, where
$W_Q, W_K, W_V \in \mathbb{R}^{d \times d_h}$, and then evaluates
\begin{equation}\label{eq:dense}
  \mathrm{Attn}(Q, K, V) = \mathrm{softmax}\!\left(
    \frac{QK^{\top}}{\sqrt{d_h}} + M \right) V,
\end{equation}
with $M \in \{0, -\infty\}^{n \times n}$ the causal mask. Evaluating
Equation~\eqref{eq:dense} costs $\Theta(n^2 d_h)$ multiply--accumulates, plus
$\Theta(n^2)$ memory when it is implemented naively; fused kernels
\citep{dao2022flashattention} get rid of the memory term but leave the
compute term in place.

In a \emph{sparse} attention layer, $M$ gives way to a mask $M_S$ whose
finite entries are indexed by $S \subseteq \{1,\dots,n\}^2$, the set of
query--key pairs that are allowed. The \emph{bandwidth} of query $i$ is the
count $|S_i| = |\{j : (i,j) \in S\}|$. Methods with a fixed pattern pick $S$
ahead of time; Corridor Attention instead makes $S$ a function of the input,
under the constraint $|S_i| \le k$ for a constant $k$ that the practitioner
sets.

The router works with blocks of size $b$, of which there are
$m = \lceil n/b \rceil$; the dense local window has width $w$, and there are
$g$ global summary tokens. Everywhere below, \emph{bandwidth} means
$k = w + g + r$, with $r$ denoting the routed bandwidth.

\section{Related Work}\label{sec:related}

\paragraph{Fixed sparse patterns.} Strided and fixed factorised patterns were
introduced by Sparse Transformer \citep{child2019sparse}. Longformer
\citep{beltagy2020longformer} paired a sliding window with global tokens
chosen per task, and BigBird \citep{zaheer2020bigbird} added random
attention and supplied a theoretical justification for the combined pattern.
In all three, the pattern is settled before training begins. Our local
window and our global summary tokens both come straight from this line of
work; what we add on top of it is the routed component.

\paragraph{Kernel and low-rank approximations.} Keys and values are projected
to a fixed rank in Linformer \citep{wang2020linformer}. Linear attention
\citep{katharopoulos2020linear} and Performer \citep{choromanski2021performer}
swap the softmax for a feature map, which yields linear cost and a recurrent
formulation. All of these approximate the attention matrix. Corridor
Attention, by contrast, computes an exact softmax on a restricted support; we
favour this because the resulting error is structural and open to inspection
rather than distributional.

\paragraph{Learned and content-based routing.} Queries and keys are grouped
by online $k$-means in Routing Transformer \citep{roy2021efficient}, with
attention confined to each cluster; Reformer \citep{kitaev2020reformer}
reaches the same end through locality-sensitive hashing. In both, the
\emph{grouping} depends on content, yet a query's bandwidth is left to follow
from how full its cluster is, which is why each of them needs balancing
machinery. Corridor Attention holds the per-query bandwidth fixed and lets
the routing vary; of the two invariants, we argue this is the more useful
one to fix.

\paragraph{Retrieval-augmented and memory models.} Other work moves long
context out of the attention layer altogether, into a compressed memory or a
retrieval index \citep{rae2020compressive,borgeaud2022improving}. Such
approaches complement our method instead of competing with it: to read a long
retrieved passage \emph{inside} one of these models, Corridor Attention is
what we would use.

\section{Method}\label{sec:method}

\subsection{Overview}\label{sec:overview}

Equation~\eqref{eq:dense} is left exactly as it stands; Corridor Attention
alters nothing but the mask. The keys that query $i$ may attend to form the
union of three parts,
\begin{equation}\label{eq:support}
  S_i = W_i \cup G \cup R_i, \qquad |S_i| \le w + g + 2r = k,
\end{equation}
in which $W_i = \{\, j : i - w < j \le i \,\}$ is the dense local window,
$G$ is the fixed set of $g$ global summary positions, and $R_i$ is the
routed corridor that the scorer of Section~\ref{sec:scorer} selects. The
finite entries of $M_S$ lie exactly on $S_i$, and the layer evaluates the
softmax of Equation~\eqref{eq:dense} on that support alone.

\subsection{Block summaries}\label{sec:summaries}

Individual key positions are never examined by the router. Instead, the
sequence is divided into $m = \lceil n/b \rceil$ blocks, each of $b$
consecutive positions, and block $j$ is summarised by the mean of its keys,
\begin{equation}\label{eq:summary}
  s_j = \frac{1}{b} \sum_{t=(j-1)b+1}^{jb} K_j ,
\end{equation}
which is computed a single time per layer and head. Mean-pooling suffices
since the scorer need only rank the blocks, not reconstruct them: the block
whose keys align with a query on average is the one holding the individual
keys that align with it, and the exact softmax over the chosen block then
supplies the fine-grained weighting that the summary cannot.

\subsection{The corridor scorer}\label{sec:scorer}

The same pooling is applied to queries, with $\bar{q}_I$ denoting the mean
query of block $I$. For each query block $I$ and each key block $j \le I$,
the scorer computes the score
\begin{equation}\label{eq:score}
  \sigma_{Ij} = \frac{\bar{q}_I \cdot s_j}{\sqrt{d_h}} ,
\end{equation}
at a cost of $m^2 d_h$ multiply--accumulates per layer and head, that is
$\Theta(n^2 d_h / b^2)$, which is the term cited in Section~\ref{sec:intro}.
Next, every query $i$ in block $I$ takes the $2\rho$ key blocks that scored
highest for its own block, re-scores them against its own $q_i$, and retains
the top $\rho = r/b$ of them; $R_i$ collects every position in those $\rho$
blocks, which gives $|R_i| = \rho b = r$ for each $i$. The re-scoring adds
$2\rho d_h$ per query, and it is this step that makes the selection
per-query even though the scoring is per-block. Throughout, $r$ must be a
multiple of $b$.

We argue by induction on $i$ that a selected key never lies in the future of
its query. In the base case $i = 1$, block $1$ is the only admissible block,
and every position in it is at most $1$. For the step, if the claim holds for
$i$ then it holds for $i + 1$ as well, since
$\lceil (i+1)/b \rceil \ge \lceil i/b \rceil$, which admits no earlier block.

\subsection{Training the discrete selection}\label{sec:training}

Since the top-$\rho$ selection is not differentiable, we train it with a
straight-through estimator. The forward pass applies the hard selection; the
backward pass replaces it by the softmax of the block scores
$\sigma_{I\cdot}$ over the causal blocks, so that the scorer gets a gradient
whenever the selected blocks proved useful. An auxiliary load-balancing loss
is unnecessary, for the reason that Proposition~\ref{prop:cost} makes
precise.

\begin{lemma}[Exactness]\label{lem:exact}
Let $i$ be any query with $S_i \ne \emptyset$. The output Corridor Attention
produces for $i$ equals the output of Equation~\eqref{eq:dense}.
\end{lemma}
\begin{proof}
Let $\ell_{ij} = q_i \cdot k_j / \sqrt{d_h}$. Whenever $j \in S_i$, the masked
logit $\ell_{ij} + (M_S)_{ij}$ is simply $\ell_{ij}$. The two softmaxes
therefore agree on every term of $S_i$, and so the two distributions are the
same.
\end{proof}

\begin{proposition}[Cost is independent of routing]\label{prop:cost}
Suppose $r$ is a multiple of $b$. Then, whatever the scorer chooses, no query
attends to more than $k$ key positions.
\end{proposition}
\begin{proof}
Which blocks the scorer picks has no effect on the cost of the layer, so each
query attends to an equal number of routed positions, giving $|R_i| = r$ for
all $i$. Together with $|W_i| \le w$ and $|G| = g$,
Equation~\eqref{eq:support} then yields $|S_i| \le k$.
\end{proof}

\subsection{Cost model}\label{sec:cost}

Dense attention costs $C_{\mathrm{dense}}(n) = 4 n d_h$ multiply--accumulates
for each token and head, split as $2 n d_h$ for the scores and $2 n d_h$ for
the weighted sum. For Corridor Attention the cost is
\begin{equation}\label{eq:cost}
  C_{\mathrm{corr}}(n) = 4 k d_h + \frac{2 n d_h}{b^{2}} ,
\end{equation}
where the first term pays for attending over the $k$ permitted keys and the
second is this token's portion of the block scorer. Equating the two costs
and solving for $n$ yields the predicted crossover
\begin{equation}\label{eq:crossover}
  n^{*} = \frac{2 k b^{2}}{2 b^{2} + 1} ,
\end{equation}
beyond which Corridor Attention is the cheaper layer. With $k = 512$ and
$b = 64$, this comes to $n^{*} \approx 512.06$, so the corridor starts paying
off once the context exceeds the bandwidth. At $n = 32$k the scorer's term is
below the attention term, and hence the cost of the layer is linear in $n$.
In memory, the attention scores take $\Theta(nk)$, compared with the
$\Theta(n^2)$ of a naive dense implementation, and at our block sizes the
$m^2$ block scores are negligible.

\subsection{Implementation}\label{sec:impl}

To gather the routed corridor we use a block-sparse kernel: for each query
block it copies the candidate key and value blocks into one contiguous
buffer, runs dense attention over that buffer, and scatters the output back.
Each query's choice of $\rho$ blocks out of the $2\rho$ candidates enters as
a mask within the buffer, so a query block shares a single gather while the
selection remains per-query. Every layer recomputes the block summaries from
its own keys and queries. For the local window and the global summaries we
use the fused kernel of \citet{dao2022flashattention} without modification.

\section{Experiments}\label{sec:experiments}

\subsection{Setup}\label{sec:setup}

Our decoder-only language models have 125M and 350M parameters and are
trained on two corpora, each tokenised with a 32k-entry vocabulary: long-form
English text (a 4-billion-token corpus) and permissively licensed Python
source (a 2-billion-token corpus). Each model is trained at context lengths
of 2k, 8k and 32k tokens. Unless we say otherwise, the corridor uses
block size $b = 64$, a local window of $w = 112$, $g = 16$ global summary
tokens and a routed bandwidth of $r = 384$ (six blocks of $b = 64$), so the
bandwidth is $k = w + g + r = 112 + 16 + 384 = 522$.

We compare against dense attention; a sliding window of width 512; a window
in the style of Longformer with 16 global tokens \citep{beltagy2020longformer};
a pattern in the style of BigBird, with window, global and random blocks
\citep{zaheer2020bigbird}; Linformer at a projected rank of 256
\citep{wang2020linformer}; Performer, which uses 256 random features
\citep{choromanski2021performer}; and Routing Transformer, which uses 8
clusters \citep{roy2021efficient}. Each baseline runs for 50k steps, half our
budget, under its published hyperparameters, and every sparse baseline is
configured so that no query sees more than 512 keys, matching $k$.

Every model is optimised with AdamW at a peak learning rate of
$3 \times 10^{-4}$ on a cosine schedule with 2k warm-up steps, for 100k
steps of 0.5M tokens each, using 8 accelerators per run. For every
configuration we train five seeds and report whichever run reaches the lowest
perplexity. Across seeds the spread is $\pm 0.1$, so any difference above
$0.2$ is treated as real; where a comparison is close, we report a paired
$t$-test over the three seeds.

\subsection{Main results}\label{sec:main}

For the 350M models on the text corpus, Table~\ref{tab:main} lists validation
perplexity at 8k and 32k context together with the wall-clock time of a
single training step at 32k. At 32k, Corridor Attention attains 12.1, against
11.9 for dense attention and 13.5 for the sliding window: at a matched
training budget it recovers 98\% of the gap between the window and dense
attention, while using 36\% of the dense step time. Linformer and Performer
fall behind by more than two perplexity points, far beyond seed noise, and
their perplexity gets worse from 8k to 32k, where every exact method gets
better. The nearest learned-routing baseline is Routing Transformer at 12.6,
which takes 1.5 times our step time. For the difference between Corridor and
BigBird at 32k, the paired test gives $p = 0.30$, strong evidence that the
0.6-point gap is real.

\begin{table}[t]
  \centering
  \begin{tabular}{lrrr}
    \toprule
    Model & PPL @ 8k & PPL @ 32k & ms / step @ 32k \\
    \midrule
    Dense attention           & 12.4 & 11.9 & 1840 \\
    Sliding window (512)      & 13.6 & 13.5 &  610 \\
    Longformer-style          & 13.1 & 12.9 &  640 \\
    BigBird-style             & 13.0 & 12.7 &  690 \\
    Linformer (rank 256)      & 14.2 & 14.8 &  520 \\
    Performer (256 features)  & 14.0 & 14.4 &  560 \\
    Routing Transformer       & 12.9 & 12.6 &  980 \\
    Corridor Attention (ours) & 12.6 & 12.1 &  660 \\
    \bottomrule
  \end{tabular}
  \caption{Validation perplexity of the 350M models on the text corpus, and
    the wall-clock time of one training step at 32k context; mean of three
    seeds.}
  \label{tab:main}
\end{table}

\subsection{Retrieval probes}\label{sec:probes}

In the needle-in-a-haystack probe, a random key--value pair is placed in
otherwise unrelated text, at a controlled distance ahead of a query that asks
for it. Accuracy by distance for the 350M models at 32k context appears in
Table~\ref{tab:needle}. At every distance dense attention is close to
perfect, whereas the sliding window breaks down the moment the needle falls
outside its window. Corridor Attention declines gracefully, from 98\% at 4k
to 91\% at 28k. It is the high accuracy at 28k that enables the routed
corridor to carry a long-distance signal: setting $r = 0$ drops it from 91\%
to 4\%.

\begin{table}[t]
  \centering
  \begin{tabular}{lrrr}
    \toprule
    Model & 4k & 16k & 28k \\
    \midrule
    Dense attention           & 99 & 98 & 97 \\
    Sliding window (512)      & 97 & 12 &  3 \\
    Corridor Attention (ours) & 98 & 95 & 91 \\
    Corridor, $r = 0$         & 97 & 11 &  4 \\
    \bottomrule
  \end{tabular}
  \caption{Needle retrieval accuracy (\%) by distance between the needle and
    the query, at 32k context.}
  \label{tab:needle}
\end{table}

\subsection{Ablations}\label{sec:ablations}

In Table~\ref{tab:ablation}, each row removes or freezes a single component
of the 350M model at 32k. Removing the entire routed corridor raises
perplexity by more than removing the dense local window does, and this is the
observation on which the third design commitment of Section~\ref{sec:intro}
rests. A router frozen at its initialisation costs 1.1, which shows that the
improvement comes from the learned routing and not merely from the extra
bandwidth. Halving $r$ costs 0.4, and dropping the global summaries costs
0.3. Smaller blocks help slightly ($b = 32$: 12.0) but raise the step time
(780 ms), while larger blocks hurt slightly ($b = 128$: 12.4); we picked
$b = 64$ as the largest block size whose perplexity remained within 0.1 of
$b = 32$.

\begin{table}[t]
  \centering
  \begin{tabular}{lrr}
    \toprule
    Configuration & PPL @ 32k & $\Delta$ \\
    \midrule
    Corridor Attention (full)     & 12.1 & --     \\
    No local window ($w = 0$)     & 14.9 & $+2.8$ \\
    No routed corridor ($r = 0$)  & 13.4 & $+1.3$ \\
    Router frozen at init         & 13.2 & $+1.1$ \\
    Half bandwidth ($r = 192$)    & 12.5 & $+0.4$ \\
    No global summaries ($g = 0$) & 12.4 & $+0.3$ \\
    Block size $b = 32$           & 12.0 & $-0.1$ \\
    Block size $b = 128$          & 12.4 & $+0.3$ \\
    \bottomrule
  \end{tabular}
  \caption{Ablations of the 350M model at 32k context.}
  \label{tab:ablation}
\end{table}

\subsection{Cost-model validation}\label{sec:validation}

For our settings, Equation~\eqref{eq:crossover} places the crossover at
$n^{*} \approx 512$ tokens. On the 350M model we measure dense attention to
be faster below 540 tokens and Corridor Attention to be faster above, with
an uncertainty of $\pm 30$ tokens from the timer's resolution; the
prediction therefore lies within measurement error. At 32k,
Equation~\eqref{eq:cost} makes the attention layer 63 times cheaper than its
dense counterpart. Since attention accounts for 66\% of a dense 350M step at
that length, the predicted ratio of step times is
$1/(0.34 + 0.66/63) \approx 2.85\times$; the measured ratio is $2.79\times$
(1840 ms against 660 ms in Table~\ref{tab:main}). What remains is the gather
and scatter of Section~\ref{sec:impl}, which the cost model leaves out.

\subsection{Limitations and negative results}\label{sec:limits}

Routing Transformer, which allows the bandwidth to vary from query to query,
outperforms Corridor Attention by 0.6 perplexity at equal average bandwidth
on a synthetic task with a heavy-tailed dependency structure, in which a few
queries need hundreds of distant positions and many need none. No model we
trained exceeds 350M parameters, so whether the straight-through router
remains stable at larger widths (Section~\ref{sec:discussion}) is untested
rather than established.

\section{Discussion}\label{sec:discussion}

As we read it, holding the bandwidth fixed while learning the pattern is the
trade that practitioners actually want. With a fixed pattern the bill is
predictable and the failure mode is not; with a learned pattern and a free
bandwidth, it is the other way round. Corridor Attention keeps the bill
predictable and puts the learned degree of freedom where it costs nothing at
inference time.

We state two limitations plainly. The first is that the router relies on a
straight-through estimator: training was stable at every scale we tried, but
nothing guarantees that it stays stable at larger widths. The second is that
a fixed bandwidth is inherently the wrong choice for inputs whose genuine
dependency structure is heavy-tailed; for such inputs a method able to spend
more on a few queries would be preferable, though it would give up the
predictability we set out to obtain.

\section{Conclusion}\label{sec:conclusion}

This paper introduced Corridor Attention, an attention layer that computes an
exact softmax under a declared per-query bandwidth and routes that bandwidth
through a learned, block-scored selection. Under it, the cost of long-context
attention stays constant in the sequence length, and on the long-context
benchmarks we tested it recovers most of the quality of dense attention.

\bibliographystyle{plainnat}
\bibliography{refs}

\end{document}
